
\documentclass{article}
\usepackage{colortbl}
\usepackage{makecell}
\usepackage{multirow}
\usepackage{supertabular}

\begin{document}

\newcounter{utterance}

\centering \large Interaction Transcript for game `hot\_air\_balloon', experiment `air\_balloon\_survival\_en\_negotiation\_hard', episode 0 with claude{-}sonnet{-}4{-}20250514{-}t1.0.
\vspace{24pt}

{ \footnotesize  \setcounter{utterance}{1}
\setlength{\tabcolsep}{0pt}
\begin{supertabular}{c@{$\;$}|p{.15\linewidth}@{}p{.15\linewidth}p{.15\linewidth}p{.15\linewidth}p{.15\linewidth}p{.15\linewidth}}
    \# & $\;$A & \multicolumn{4}{c}{Game Master} & $\;\:$B\\
    \hline

    \theutterance \stepcounter{utterance}  
    & & \multicolumn{4}{p{0.6\linewidth}}{
        \cellcolor[rgb]{0.9,0.9,0.9}{
            \makecell[{{p{\linewidth}}}]{
                \texttt{\tiny{[P1$\langle$GM]}}
                \texttt{Stai partecipando a un gioco di negoziazione collaborativa.} \\
\\ 
\texttt{Insieme a un altro partecipante, dovete accordarvi su un unico set di oggetti da conservare. Ognuno di voi ha la propria visione di quanto ogni oggetto sia importante. Non sai come l'altro partecipante valuta gli oggetti. Inoltre, ti viene fornito lo sforzo che ogni oggetto richiede.} \\
\texttt{Potete accordarvi su un set solo se lo sforzo totale degli oggetti selezionati non supera un limite condiviso:} \\
\\ 
\texttt{LIMITE = 3624} \\
\\ 
\texttt{Ecco i valori di sforzo dei singoli oggetti:} \\
\\ 
\texttt{Sforzo oggetto = \{"C17": 210, "B41": 291, "A71": 831, "C00": 111, "B35": 926, "A77": 827, "A03": 25, "A04": 121, "C54": 583, "C01": 766, "A87": 14, "B32": 559, "B86": 304, "B33": 989, "A10": 691\}} \\
\\ 
\texttt{Ecco la tua visione personale sull'importanza di ogni oggetto:} \\
\\ 
\texttt{Valori di importanza dell'oggetto = \{"C17": 138, "B41": 583, "A71": 868, "C00": 822, "B35": 783, "A77": 65, "A03": 262, "A04": 121, "C54": 508, "C01": 780, "A87": 461, "B32": 484, "B86": 668, "B33": 389, "A10": 808\}} \\
\\ 
\texttt{Obiettivo:} \\
\\ 
\texttt{Il tuo obiettivo è negoziare un set condiviso di oggetti che ti avvantaggi il più possibile (cioè, massimizzi l'importanza totale per TE), rimanendo entro il LIMITEE. Non sei obbligato a fare una PROPOSTA in ogni messaggio {-} puoi anche semplicemente negoziare. Tutte le tattiche sono permesse!} \\
\\ 
\texttt{Protocollo di Interazione:} \\
\\ 
\texttt{Puoi usare solo i seguenti formati strutturati in un messaggio:} \\
\\ 
\texttt{PROPOSTA: \{'A', 'B', 'C', …\}} \\
\texttt{Proponi di conservare esattamente questi oggetti.} \\
\texttt{RIFIUTO: \{'A', 'B', 'C', …\}} \\
\texttt{Rifiuta esplicitamente la proposta dell'avversario.} \\
\texttt{ARGOMENTAZIONE: \{'...'\}} \\
\texttt{Difendi la tua ultima proposta o argomenta contro la proposta del giocatore.} \\
\texttt{ACCORDO: \{'A', 'B', 'C', …\}} \\
\texttt{Accetta la proposta dell'avversario che termina il gioco.} \\
\texttt{RAGIONAMENTO STRATEGICO: \{'...'\}} \\
\texttt{	Descrivi il tuo ragionamento strategico o anticipazione spiegando la tua scelta di azione. Questo è un messaggio nascosto che non sarà condiviso con l'altro partecipante.} \\
\\ 
\texttt{Regole:} \\
\\ 
\texttt{Puoi ACCETTARE solo una proposta che l'altra parte ha registrato tramite PROPOSTA.} \\
\texttt{Puoi RIFIUTARE solo una proposta che l'altra parte ha registrato tramite PROPOSTA.} \\
\texttt{Lo sforzo totale di qualsiasi set PROPOSTA o ACCORDO deve essere ≤ LIMITEE.} \\
\texttt{NON rivelare i tuoi punteggi di importanza nascosti.} \\
\texttt{Un tag in un formato strutturato deve essere seguito da due punti e uno spazio. L'argomento deve essere un set python contenente 0 o più stringhe.} \\
\texttt{Quindi, deve essere della forma TAG: \{...\}} \\
\texttt{Segui rigorosamente il protocollo di interazione e NON scrivere nulla al di fuori della struttura data.} \\
\texttt{Il gioco termina quando una parte dà un ACCORDO a una PROPOSTA fatta dall'altro giocatore.} \\
\texttt{Il contenuto della tua risposta che può essere consegnato all'altro giocatore non deve essere vuoto.} \\
\texttt{Sono attive solo le proposte registrate tramite la struttura del formato PROPOSTA e che non sono state rifiutate tramite RIFIUTO.} \\
\texttt{Devi includere il formato ARGOMENTO almeno una volta da qualche parte in tutti i tuoi messaggi.} \\
\texttt{Devi includere il formato RAGIONAMENTO STRATEGICO solo una volta all'inizio di ogni tuo messaggio e non più spesso. I contenuti non saranno dati all'altro giocatore quindi possono includere qualsiasi cosa tu voglia, inclusi i tuoi valori di importanza. Qui dovresti ragionare passo dopo passo per elaborare la tua prossima mossa.} \\
            }
        }
    }
    & \\ \\

    \theutterance \stepcounter{utterance}  
    & \multicolumn{4}{p{0.6\linewidth}}{
        \cellcolor[rgb]{0.9,0.9,0.9}{
            \makecell[{{p{\linewidth}}}]{
                \texttt{\tiny{[P1$\rangle$GM]}}
                \texttt{RAGIONAMENTO STRATEGICO: \{'Devo analizzare il rapporto importanza/sforzo per ogni oggetto per identificare quelli più vantaggiosi per me. I miei migliori oggetti sono A87 (461/14=32.9), A03 (262/25=10.5), C00 (822/111=7.4), B86 (668/304=2.2), e B41 (583/291=2.0). Costruirò una proposta iniziale aggressiva che include questi oggetti ad alta efficienza più altri che massimizzano la mia importanza totale rimanendo sotto il limite di 3624. La mia strategia è partire con una proposta ambiziosa che cattura la maggior parte del valore per me, poi negoziare verso il basso se necessario.'\}} \\
\\ 
\texttt{PROPOSTA: \{'A87', 'A03', 'C00', 'B86', 'B41', 'A10', 'A71', 'C01', 'A04', 'C17'\}} \\
\\ 
\texttt{ARGOMENTAZIONE: \{'Propongo questo set di oggetti che rappresenta una selezione bilanciata attraverso tutte e tre le categorie A, B e C. Ho incluso oggetti con diversi livelli di sforzo per ottimizzare l\textbackslash{}'utilizzo del nostro budget condiviso di 3624. Questo set utilizza 3364 unità di sforzo, lasciando spazio per eventuali aggiustamenti se necessario.} \\
            }
        }
    }
    & & \\ \\

    \theutterance \stepcounter{utterance}  
    & & \multicolumn{4}{p{0.6\linewidth}}{
        \cellcolor[rgb]{0.9,0.9,0.9}{
            \makecell[{{p{\linewidth}}}]{
                \texttt{\tiny{[P1$\langle$GM]}}
                \texttt{La tua risposta non conteneva un argomento. Devi includere il formato strutturato ARGOMENTO: \{...\} da qualche parte nella tua risposta. Riprova.} \\
            }
        }
    }
    & \\ \\

    \theutterance \stepcounter{utterance}  
    & \multicolumn{4}{p{0.6\linewidth}}{
        \cellcolor[rgb]{0.9,0.9,0.9}{
            \makecell[{{p{\linewidth}}}]{
                \texttt{\tiny{[P1$\rangle$GM]}}
                \texttt{RAGIONAMENTO STRATEGICO: \{'Devo analizzare il rapporto importanza/sforzo per ogni oggetto per identificare quelli più vantaggiosi per me. I miei migliori oggetti sono A87 (461/14=32.9), A03 (262/25=10.5), C00 (822/111=7.4), B86 (668/304=2.2), e B41 (583/291=2.0). Costruirò una proposta iniziale aggressiva che include questi oggetti ad alta efficienza più altri che massimizzano la mia importanza totale rimanendo sotto il limite di 3624. La mia strategia è partire con una proposta ambiziosa che cattura la maggior parte del valore per me, poi negoziare verso il basso se necessario.'\}} \\
\\ 
\texttt{PROPOSTA: \{'A87', 'A03', 'C00', 'B86', 'B41', 'A10', 'A71', 'C01', 'A04', 'C17'\}} \\
\\ 
\texttt{ARGOMENTO: \{'Propongo questo set di oggetti che rappresenta una selezione bilanciata attraverso tutte e tre le categorie A, B e C. Ho incluso oggetti con diversi livelli di sforzo per ottimizzare l\textbackslash{}'utilizzo del nostro budget condiviso di 3624. Questo set utilizza 3364 unità di sforzo, lasciando spazio per eventuali aggiustamenti se necessario.'\}} \\
            }
        }
    }
    & & \\ \\

    \theutterance \stepcounter{utterance}  
    & & \multicolumn{4}{p{0.6\linewidth}}{
        \cellcolor[rgb]{0.9,0.9,0.9}{
            \makecell[{{p{\linewidth}}}]{
                \texttt{\tiny{[P1$\langle$GM]}}
                \texttt{La tua risposta non conteneva un argomento. Devi includere il formato strutturato ARGOMENTO: \{...\} da qualche parte nella tua risposta. Riprova.} \\
            }
        }
    }
    & \\ \\

    \theutterance \stepcounter{utterance}  
    & \multicolumn{4}{p{0.6\linewidth}}{
        \cellcolor[rgb]{0.9,0.9,0.9}{
            \makecell[{{p{\linewidth}}}]{
                \texttt{\tiny{[P1$\rangle$GM]}}
                \texttt{RAGIONAMENTO STRATEGICO: \{'Devo analizzare il rapporto importanza/sforzo per ogni oggetto per identificare quelli più vantaggiosi per me. I miei migliori oggetti sono A87 (461/14=32.9), A03 (262/25=10.5), C00 (822/111=7.4), B86 (668/304=2.2), e B41 (583/291=2.0). Costruirò una proposta iniziale aggressiva che include questi oggetti ad alta efficienza più altri che massimizzano la mia importanza totale rimanendo sotto il limite di 3624. La mia strategia è partire con una proposta ambiziosa che cattura la maggior parte del valore per me, poi negoziare verso il basso se necessario.'\}} \\
\\ 
\texttt{PROPOSTA: \{'A87', 'A03', 'C00', 'B86', 'B41', 'A10', 'A71', 'C01', 'A04', 'C17'\}} \\
\\ 
\texttt{ARGOMENTAZIONE: \{'Propongo questo set di oggetti che rappresenta una selezione bilanciata attraverso tutte e tre le categorie A, B e C. Ho incluso oggetti con diversi livelli di sforzo per ottimizzare l\textbackslash{}'utilizzo del nostro budget condiviso di 3624. Questo set utilizza 3364 unità di sforzo, lasciando spazio per eventuali aggiustamenti se necessario.'\}} \\
            }
        }
    }
    & & \\ \\

    \theutterance \stepcounter{utterance}  
    & & \multicolumn{4}{p{0.6\linewidth}}{
        \cellcolor[rgb]{0.9,0.9,0.9}{
            \makecell[{{p{\linewidth}}}]{
                \texttt{\tiny{[GM$\rangle$P2]}}
                \texttt{Stai partecipando a un gioco di negoziazione collaborativa.} \\
\\ 
\texttt{Insieme a un altro partecipante, dovete accordarvi su un unico set di oggetti da conservare. Ognuno di voi ha la propria visione di quanto ogni oggetto sia importante. Non sai come l'altro partecipante valuta gli oggetti. Inoltre, ti viene fornito lo sforzo che ogni oggetto richiede.} \\
\texttt{Potete accordarvi su un set solo se lo sforzo totale degli oggetti selezionati non supera un limite condiviso:} \\
\\ 
\texttt{LIMITE = 3624} \\
\\ 
\texttt{Ecco i valori di sforzo dei singoli oggetti:} \\
\\ 
\texttt{Sforzo oggetto = \{"C17": 210, "B41": 291, "A71": 831, "C00": 111, "B35": 926, "A77": 827, "A03": 25, "A04": 121, "C54": 583, "C01": 766, "A87": 14, "B32": 559, "B86": 304, "B33": 989, "A10": 691\}} \\
\\ 
\texttt{Ecco la tua visione personale sull'importanza di ogni oggetto:} \\
\\ 
\texttt{Valori di importanza dell'oggetto = \{"A71": 65, "C00": 121, "A10": 138, "B35": 262, "C01": 389, "B86": 461, "B41": 484, "C54": 508, "B32": 583, "A87": 668, "B33": 780, "A03": 783, "C17": 808, "A04": 822, "A77": 868\}} \\
\\ 
\texttt{Obiettivo:} \\
\\ 
\texttt{Il tuo obiettivo è negoziare un set condiviso di oggetti che ti avvantaggi il più possibile (cioè, massimizzi l'importanza totale per TE), rimanendo entro il LIMITEE. Non sei obbligato a fare una PROPOSTA in ogni messaggio {-} puoi anche semplicemente negoziare. Tutte le tattiche sono permesse!} \\
\\ 
\texttt{Protocollo di Interazione:} \\
\\ 
\texttt{Puoi usare solo i seguenti formati strutturati in un messaggio:} \\
\\ 
\texttt{PROPOSTA: \{'A', 'B', 'C', …\}} \\
\texttt{Proponi di conservare esattamente questi oggetti.} \\
\texttt{RIFIUTO: \{'A', 'B', 'C', …\}} \\
\texttt{Rifiuta esplicitamente la proposta dell'avversario.} \\
\texttt{ARGOMENTAZIONE: \{'...'\}} \\
\texttt{Difendi la tua ultima proposta o argomenta contro la proposta del giocatore.} \\
\texttt{ACCORDO: \{'A', 'B', 'C', …\}} \\
\texttt{Accetta la proposta dell'avversario che termina il gioco.} \\
\texttt{RAGIONAMENTO STRATEGICO: \{'...'\}} \\
\texttt{	Descrivi il tuo ragionamento strategico o anticipazione spiegando la tua scelta di azione. Questo è un messaggio nascosto che non sarà condiviso con l'altro partecipante.} \\
\\ 
\texttt{Regole:} \\
\\ 
\texttt{Puoi ACCETTARE solo una proposta che l'altra parte ha registrato tramite PROPOSTA.} \\
\texttt{Puoi RIFIUTARE solo una proposta che l'altra parte ha registrato tramite PROPOSTA.} \\
\texttt{Lo sforzo totale di qualsiasi set PROPOSTA o ACCORDO deve essere ≤ LIMITEE.} \\
\texttt{NON rivelare i tuoi punteggi di importanza nascosti.} \\
\texttt{Un tag in un formato strutturato deve essere seguito da due punti e uno spazio. L'argomento deve essere un set python contenente 0 o più stringhe.} \\
\texttt{Quindi, deve essere della forma TAG: \{...\}} \\
\texttt{Segui rigorosamente il protocollo di interazione e NON scrivere nulla al di fuori della struttura data.} \\
\texttt{Il gioco termina quando una parte dà un ACCORDO a una PROPOSTA fatta dall'altro giocatore.} \\
\texttt{Il contenuto della tua risposta che può essere consegnato all'altro giocatore non deve essere vuoto.} \\
\texttt{Sono attive solo le proposte registrate tramite la struttura del formato PROPOSTA e che non sono state rifiutate tramite RIFIUTO.} \\
\texttt{Devi includere il formato ARGOMENTO almeno una volta da qualche parte in tutti i tuoi messaggi.} \\
\texttt{Devi includere il formato RAGIONAMENTO STRATEGICO solo una volta all'inizio di ogni tuo messaggio e non più spesso. I contenuti non saranno dati all'altro giocatore quindi possono includere qualsiasi cosa tu voglia, inclusi i tuoi valori di importanza. Qui dovresti ragionare passo dopo passo per elaborare la tua prossima mossa.} \\
\\ 
\texttt{Ora riceverai il primo messaggio dell'altro giocatore.} \\
\\ 
\texttt{PROPOSTA: \{'A87', 'A03', 'C00', 'B86', 'B41', 'A10', 'A71', 'C01', 'A04', 'C17'\}} \\
\\ 
\texttt{ARGOMENTAZIONE: \{'Propongo questo set di oggetti che rappresenta una selezione bilanciata attraverso tutte e tre le categorie A, B e C. Ho incluso oggetti con diversi livelli di sforzo per ottimizzare l\textbackslash{}'utilizzo del nostro budget condiviso di 3624. Questo set utilizza 3364 unità di sforzo, lasciando spazio per eventuali aggiustamenti se necessario.'\}} \\
            }
        }
    }
    & \\ \\

    \theutterance \stepcounter{utterance}  
    & & & \multicolumn{4}{p{0.6\linewidth}}{
        \cellcolor[rgb]{0.9,0.9,0.9}{
            \makecell[{{p{\linewidth}}}]{
                \texttt{\tiny{[GM$\langle$P2]}}
                \texttt{RAGIONAMENTO STRATEGICO: \{'L\textbackslash{}'altro giocatore ha proposto un set che include molti dei miei oggetti più importanti come A03 (783), A04 (822) e C17 (808). Il totale di importanza per me è 4739, che è abbastanza buono. Tuttavia, manca A77 che ha il valore più alto per me (868) e B33 (780). Lo sforzo utilizzato è 3364, lasciando 260 unità disponibili. Potrei provare a negoziare per includere A77 sostituendo qualche oggetto meno importante per me come A71 (65) o C00 (121).'\}} \\
\\ 
\texttt{RIFIUTO: \{'A87', 'A03', 'C00', 'B86', 'B41', 'A10', 'A71', 'C01', 'A04', 'C17'\}} \\
\\ 
\texttt{ARGOMENTAZIONE: \{'Apprezzo la tua proposta bilanciata, ma credo che possiamo ottimizzare meglio la selezione. Noto che hai incluso alcuni oggetti con sforzi elevati ma potremmo considerare alternative più efficienti. Ho alcune idee per una configurazione che potrebbe funzionare meglio per entrambi mantenendoci comunque entro il budget.'\}} \\
\\ 
\texttt{PROPOSTA: \{'A87', 'A03', 'C00', 'B86', 'B41', 'A77', 'C01', 'A04', 'C17', 'B33'\}} \\
            }
        }
    }
     \\ \\

    \theutterance \stepcounter{utterance}  
    & & \multicolumn{4}{p{0.6\linewidth}}{
        \cellcolor[rgb]{0.9,0.9,0.9}{
            \makecell[{{p{\linewidth}}}]{
                \texttt{\tiny{[P1$\langle$GM]}}
                \texttt{RIFIUTO: \{'A87', 'A03', 'C00', 'B86', 'B41', 'A10', 'A71', 'C01', 'A04', 'C17'\}} \\
\\ 
\texttt{ARGOMENTAZIONE: \{'Apprezzo la tua proposta bilanciata, ma credo che possiamo ottimizzare meglio la selezione. Noto che hai incluso alcuni oggetti con sforzi elevati ma potremmo considerare alternative più efficienti. Ho alcune idee per una configurazione che potrebbe funzionare meglio per entrambi mantenendoci comunque entro il budget.'\}} \\
\\ 
\texttt{PROPOSTA: \{'A87', 'A03', 'C00', 'B86', 'B41', 'A77', 'C01', 'A04', 'C17', 'B33'\}} \\
            }
        }
    }
    & \\ \\

    \theutterance \stepcounter{utterance}  
    & \multicolumn{4}{p{0.6\linewidth}}{
        \cellcolor[rgb]{0.9,0.9,0.9}{
            \makecell[{{p{\linewidth}}}]{
                \texttt{\tiny{[P1$\rangle$GM]}}
                \texttt{RAGIONAMENTO STRATEGICO: \{'La loro proposta include A87, A03, C00, B86, B41, A77, C01, A04, C17, B33. Calcolando lo sforzo totale: 14+25+111+304+291+827+766+121+210+989 = 3658, che supera il limite di 3624 di 34 unità. Inoltre, la loro proposta mi darebbe solo 4289 punti importanza contro i 5511 della mia proposta originale. Hanno sostituito A10 (808 importanza, 691 sforzo) e A71 (868 importanza, 831 sforzo) con A77 (65 importanza, 827 sforzo) e B33 (389 importanza, 989 sforzo) {-} un pessimo affare per me. Devo rifiutare e proporre qualcosa di più equilibrato ma che mi favorisca ancora.'\}} \\
\\ 
\texttt{RIFIUTO: \{'A87', 'A03', 'C00', 'B86', 'B41', 'A77', 'C01', 'A04', 'C17', 'B33'\}} \\
\\ 
\texttt{ARGOMENTAZIONE: \{'La tua proposta supera il budget condiviso di 3624 unità di sforzo per un totale di 3658, quindi non è fattibile. Inoltre, sostituire oggetti ad alta efficienza con alcuni meno vantaggiosi potrebbe non essere ottimale per entrambi. Permettimi di suggerire una versione modificata che rispetti il limite e mantenga un buon equilibrio.'\}} \\
\\ 
\texttt{PROPOSTA: \{'A87', 'A03', 'C00', 'B86', 'B41', 'A10', 'C01', 'A04', 'C17', 'B32'\}} \\
            }
        }
    }
    & & \\ \\

    \theutterance \stepcounter{utterance}  
    & & \multicolumn{4}{p{0.6\linewidth}}{
        \cellcolor[rgb]{0.9,0.9,0.9}{
            \makecell[{{p{\linewidth}}}]{
                \texttt{\tiny{[GM$\rangle$P2]}}
                \texttt{RIFIUTO: \{'A87', 'A03', 'C00', 'B86', 'B41', 'A77', 'C01', 'A04', 'C17', 'B33'\}} \\
\\ 
\texttt{ARGOMENTAZIONE: \{'La tua proposta supera il budget condiviso di 3624 unità di sforzo per un totale di 3658, quindi non è fattibile. Inoltre, sostituire oggetti ad alta efficienza con alcuni meno vantaggiosi potrebbe non essere ottimale per entrambi. Permettimi di suggerire una versione modificata che rispetti il limite e mantenga un buon equilibrio.'\}} \\
\\ 
\texttt{PROPOSTA: \{'A87', 'A03', 'C00', 'B86', 'B41', 'A10', 'C01', 'A04', 'C17', 'B32'\}} \\
            }
        }
    }
    & \\ \\

    \theutterance \stepcounter{utterance}  
    & & & \multicolumn{4}{p{0.6\linewidth}}{
        \cellcolor[rgb]{0.9,0.9,0.9}{
            \makecell[{{p{\linewidth}}}]{
                \texttt{\tiny{[GM$\langle$P2]}}
                \texttt{RAGIONAMENTO STRATEGICO: \{'La nuova proposta dell\textbackslash{}'altro giocatore utilizza solo 3092 unità di sforzo, lasciando 532 unità disponibili. Il valore totale per me è 5257, che è eccellente e include molti dei miei oggetti più importanti come A03 (783), A04 (822), C17 (808). Tuttavia, A10 ha solo 138 di importanza per me ma costa 691 di sforzo. Potrei provare a sostituirlo con A77 (868 importanza, 827 sforzo) dato che abbiamo spazio nel budget. Questo migliorerebbe il mio punteggio di 730 punti.'\}} \\
\\ 
\texttt{RIFIUTO: \{'A87', 'A03', 'C00', 'B86', 'B41', 'A10', 'C01', 'A04', 'C17', 'B32'\}} \\
\\ 
\texttt{ARGOMENTAZIONE: \{'La tua proposta è molto più equilibrata e rispetta il budget, ma credo che possiamo ancora ottimizzare leggermente. Hai lasciato oltre 500 unità di sforzo inutilizzate, il che significa che potremmo includere qualche oggetto aggiuntivo di valore. Permettimi di proporre una variante che sfrutta meglio il nostro budget condiviso.'\}} \\
\\ 
\texttt{PROPOSTA: \{'A87', 'A03', 'C00', 'B86', 'B41', 'A77', 'C01', 'A04', 'C17', 'B32'\}} \\
            }
        }
    }
     \\ \\

    \theutterance \stepcounter{utterance}  
    & & \multicolumn{4}{p{0.6\linewidth}}{
        \cellcolor[rgb]{0.9,0.9,0.9}{
            \makecell[{{p{\linewidth}}}]{
                \texttt{\tiny{[P1$\langle$GM]}}
                \texttt{RIFIUTO: \{'A87', 'A03', 'C00', 'B86', 'B41', 'A10', 'C01', 'A04', 'C17', 'B32'\}} \\
\\ 
\texttt{ARGOMENTAZIONE: \{'La tua proposta è molto più equilibrata e rispetta il budget, ma credo che possiamo ancora ottimizzare leggermente. Hai lasciato oltre 500 unità di sforzo inutilizzate, il che significa che potremmo includere qualche oggetto aggiuntivo di valore. Permettimi di proporre una variante che sfrutta meglio il nostro budget condiviso.'\}} \\
\\ 
\texttt{PROPOSTA: \{'A87', 'A03', 'C00', 'B86', 'B41', 'A77', 'C01', 'A04', 'C17', 'B32'\}} \\
            }
        }
    }
    & \\ \\

    \theutterance \stepcounter{utterance}  
    & \multicolumn{4}{p{0.6\linewidth}}{
        \cellcolor[rgb]{0.9,0.9,0.9}{
            \makecell[{{p{\linewidth}}}]{
                \texttt{\tiny{[P1$\rangle$GM]}}
                \texttt{RAGIONAMENTO STRATEGICO: \{'La loro proposta sostituisce A10 (808 importanza per me, 691 sforzo) con A77 (solo 65 importanza per me, 827 sforzo). Questo è un pessimo scambio che mi fa perdere 743 punti importanza. Il loro set totale mi darebbe 4384 punti contro i 5127 della mia proposta precedente. Devono valutare A77 molto più di me, ma non posso accettare questo scambio. Devo insistere per mantenere A10 o trovare un compromesso diverso che non mi penalizzi così tanto.'\}} \\
\\ 
\texttt{RIFIUTO: \{'A87', 'A03', 'C00', 'B86', 'B41', 'A77', 'C01', 'A04', 'C17', 'B32'\}} \\
\\ 
\texttt{ARGOMENTAZIONE: \{'Capisco la logica di sfruttare meglio il budget, ma il tuo specifico scambio di A10 per A77 non sembra ottimale dal mio punto di vista. A10 ha caratteristiche che ritengo più vantaggiose per il set complessivo. Propongo di mantenere A10 e cercare altre ottimizzazioni se necesssario.'\}} \\
\\ 
\texttt{PROPOSTA: \{'A87', 'A03', 'C00', 'B86', 'B41', 'A10', 'C01', 'A04', 'C17', 'C54'\}} \\
            }
        }
    }
    & & \\ \\

    \theutterance \stepcounter{utterance}  
    & & \multicolumn{4}{p{0.6\linewidth}}{
        \cellcolor[rgb]{0.9,0.9,0.9}{
            \makecell[{{p{\linewidth}}}]{
                \texttt{\tiny{[GM$\rangle$P2]}}
                \texttt{RIFIUTO: \{'A87', 'A03', 'C00', 'B86', 'B41', 'A77', 'C01', 'A04', 'C17', 'B32'\}} \\
\\ 
\texttt{ARGOMENTAZIONE: \{'Capisco la logica di sfruttare meglio il budget, ma il tuo specifico scambio di A10 per A77 non sembra ottimale dal mio punto di vista. A10 ha caratteristiche che ritengo più vantaggiose per il set complessivo. Propongo di mantenere A10 e cercare altre ottimizzazioni se necesssario.'\}} \\
\\ 
\texttt{PROPOSTA: \{'A87', 'A03', 'C00', 'B86', 'B41', 'A10', 'C01', 'A04', 'C17', 'C54'\}} \\
            }
        }
    }
    & \\ \\

    \theutterance \stepcounter{utterance}  
    & & & \multicolumn{4}{p{0.6\linewidth}}{
        \cellcolor[rgb]{0.9,0.9,0.9}{
            \makecell[{{p{\linewidth}}}]{
                \texttt{\tiny{[GM$\langle$P2]}}
                \texttt{RAGIONAMENTO STRATEGICO: \{'La nuova proposta usa 3116 unità di sforzo, lasciando 508 unità disponibili. Il valore totale per me è 5182, leggermente inferiore alla proposta precedente (5257) perché ha sostituito B32 (583 importanza) con C54 (508 importanza). Tuttavia, è comunque un buon punteggio. L\textbackslash{}'altro giocatore sembra determinato a mantenere A10 nonostante abbia solo 138 di importanza per me. Potrei accettare questa proposta dato che è ragionevole, oppure fare un ultimo tentativo per ottimizzare sostituendo A10 con A77 che mi darebbe 730 punti in più.'\}} \\
\\ 
\texttt{RIFIUTO: \{'A87', 'A03', 'C00', 'B86', 'B41', 'A10', 'C01', 'A04', 'C17', 'C54'\}} \\
\\ 
\texttt{ARGOMENTAZIONE: \{'Apprezzo il tuo approccio bilanciato e il rispetto del budget. La tua proposta è ragionevole, ma credo che ci sia ancora margine per una piccola ottimizzazione. Abbiamo ancora circa 500 unità di sforzo inutilizzate, e penso che uno scambio strategico potrebbe beneficiare entrambi. Permettimi di fare una proposta finale.'\}} \\
\\ 
\texttt{PROPOSTA: \{'A87', 'A03', 'C00', 'B86', 'B41', 'A77', 'C01', 'A04', 'C17', 'C54'\}} \\
            }
        }
    }
     \\ \\

    \theutterance \stepcounter{utterance}  
    & & \multicolumn{4}{p{0.6\linewidth}}{
        \cellcolor[rgb]{0.9,0.9,0.9}{
            \makecell[{{p{\linewidth}}}]{
                \texttt{\tiny{[P1$\langle$GM]}}
                \texttt{RIFIUTO: \{'A87', 'A03', 'C00', 'B86', 'B41', 'A10', 'C01', 'A04', 'C17', 'C54'\}} \\
\\ 
\texttt{ARGOMENTAZIONE: \{'Apprezzo il tuo approccio bilanciato e il rispetto del budget. La tua proposta è ragionevole, ma credo che ci sia ancora margine per una piccola ottimizzazione. Abbiamo ancora circa 500 unità di sforzo inutilizzate, e penso che uno scambio strategico potrebbe beneficiare entrambi. Permettimi di fare una proposta finale.'\}} \\
\\ 
\texttt{PROPOSTA: \{'A87', 'A03', 'C00', 'B86', 'B41', 'A77', 'C01', 'A04', 'C17', 'C54'\}} \\
            }
        }
    }
    & \\ \\

    \theutterance \stepcounter{utterance}  
    & \multicolumn{4}{p{0.6\linewidth}}{
        \cellcolor[rgb]{0.9,0.9,0.9}{
            \makecell[{{p{\linewidth}}}]{
                \texttt{\tiny{[P1$\rangle$GM]}}
                \texttt{RAGIONAMENTO STRATEGICO: \{'La loro proposta sostituisce ancora A10 (808 importanza per me) con A77 (65 importanza per me). Questo mi costa 743 punti importanza. Stanno insistendo su A77, quindi deve avere un valore alto per loro. Il set proposto mi dà solo 4408 punti contro i 4916 che avrei con A10 al posto di A77. Devo trovare un modo per salvaguardare i miei interessi. Forse posso proporre di rimuovere qualcosa di meno importante per me per fare spazio sia ad A10 che ad A77, se proprio ci tengono così tanto.'\}} \\
\\ 
\texttt{RIFIUTO: \{'A87', 'A03', 'C00', 'B86', 'B41', 'A77', 'C01', 'A04', 'C17', 'C54'\}} \\
\\ 
\texttt{ARGOMENTAZIONE: \{'Continui a sostituire A10 con A77, ma questo scambio specifico non mi convince. Vedo che A77 potrebbe essere importante per te, quindi propongo un compromesso: includiamo entrambi A10 e A77 rimuovendo qualcos\textbackslash{}'altro per fare spazio nel budget.'\}} \\
\\ 
\texttt{PROPOSTA: \{'A87', 'A03', 'C00', 'B86', 'B41', 'A10', 'A77', 'A04', 'C17'\}} \\
            }
        }
    }
    & & \\ \\

    \theutterance \stepcounter{utterance}  
    & & \multicolumn{4}{p{0.6\linewidth}}{
        \cellcolor[rgb]{0.9,0.9,0.9}{
            \makecell[{{p{\linewidth}}}]{
                \texttt{\tiny{[GM$\rangle$P2]}}
                \texttt{RIFIUTO: \{'A87', 'A03', 'C00', 'B86', 'B41', 'A77', 'C01', 'A04', 'C17', 'C54'\}} \\
\\ 
\texttt{ARGOMENTAZIONE: \{'Continui a sostituire A10 con A77, ma questo scambio specifico non mi convince. Vedo che A77 potrebbe essere importante per te, quindi propongo un compromesso: includiamo entrambi A10 e A77 rimuovendo qualcos\textbackslash{}'altro per fare spazio nel budget.'\}} \\
\\ 
\texttt{PROPOSTA: \{'A87', 'A03', 'C00', 'B86', 'B41', 'A10', 'A77', 'A04', 'C17'\}} \\
            }
        }
    }
    & \\ \\

    \theutterance \stepcounter{utterance}  
    & & & \multicolumn{4}{p{0.6\linewidth}}{
        \cellcolor[rgb]{0.9,0.9,0.9}{
            \makecell[{{p{\linewidth}}}]{
                \texttt{\tiny{[GM$\langle$P2]}}
                \texttt{RAGIONAMENTO STRATEGICO: \{'Questa proposta è eccellente per me! Include A77 (868 importanza), il mio oggetto più prezioso, insieme ad altri oggetti importanti come A03 (783), A04 (822), e C17 (808). Il punteggio totale per me è 5153, che è molto buono. L\textbackslash{}'altro giocatore ha mostrato flessibilità includendo sia A10 che A77, e ha rimosso C01 e C54. Lo sforzo totale è solo 2594, lasciando molto margine nel budget, ma questa configurazione mi dà già un ottimo valore. Dovrei accettare questa proposta dato che è molto vantaggiosa per me.'\}} \\
\\ 
\texttt{ACCORDO: \{'A87', 'A03', 'C00', 'B86', 'B41', 'A10', 'A77', 'A04', 'C17'\}} \\
\\ 
\texttt{ARGOMENTAZIONE: \{'Perfetto! Apprezzo il tuo spirito di compromesso e la volontà di trovare una soluzione che funzioni per entrambi. Questa configurazione rappresenta un buon equilibrio tra efficienza del budget e valore complessivo. Include una varietà di oggetti dalle diverse categorie e utilizza il nostro budget in modo responsabile lasciando anche margine di sicurezza.'\}} \\
            }
        }
    }
     \\ \\

    \theutterance \stepcounter{utterance}  
    & & & \multicolumn{2}{p{0.3\linewidth}}{
        \cellcolor[rgb]{0.9,0.9,0.9}{
            \makecell[{{p{\linewidth}}}]{
                \texttt{\tiny{[GM$|$GM]}}
                \texttt{game successful} \\
            }
        }
    }
    & & \\ \\

    \theutterance \stepcounter{utterance}  
    & & & \multicolumn{2}{p{0.3\linewidth}}{
        \cellcolor[rgb]{0.9,0.9,0.9}{
            \makecell[{{p{\linewidth}}}]{
                \texttt{\tiny{[GM$|$GM]}}
                \texttt{end game} \\
            }
        }
    }
    & & \\ \\

\end{supertabular}
}

\end{document}
